\documentclass[11pt,letterpaper]{article}
\usepackage[T1]{fontenc}
\usepackage{lmodern,microtype}
\usepackage[margin=0.78in,headheight=15pt,headsep=15pt]{geometry}
\usepackage{mathtools,amssymb}
\usepackage{xcolor,booktabs,array}
\usepackage[most]{tcolorbox}
\usepackage{enumitem,fancyhdr,needspace,hyperref}

\definecolor{ink}{HTML}{162332}
\definecolor{navy}{HTML}{17324D}
\definecolor{teal}{HTML}{167D82}
\definecolor{mist}{HTML}{EDF5F4}
\definecolor{muted}{HTML}{5D6B78}
\color{ink}
\hypersetup{colorlinks=true,linkcolor=teal,urlcolor=teal,
  pdfauthor={Igor Rivin},
  pdftitle={Number Theory: Week 2 Problem Set},
  pdfsubject={Continued fractions, the Euclidean algorithm, unique factorization, and the sieve of Eratosthenes}}
\setlength{\parindent}{0pt}
\setlength{\parskip}{4pt}
\setlist[enumerate]{label=(\alph*),leftmargin=1.8em,itemsep=3pt,topsep=4pt}
\setlength{\abovedisplayskip}{7pt}
\setlength{\belowdisplayskip}{7pt}
\newcommand{\Z}{\mathbb Z}
\newcommand{\pts}[1]{\hfill\textcolor{muted}{\small[#1 points]}}
\newcommand{\partpts}[1]{\textcolor{muted}{\small(#1)}}
\newcommand{\edition}{\ifsolutions Instructor solutions\else Student handout\fi}
\newcounter{problem}
\newcommand{\problem}[2]{%
  \par\addvspace{11pt}\Needspace{6\baselineskip}%
  \stepcounter{problem}{\bfseries\color{navy}\theproblem.\ #1}\pts{#2}\par\nobreak}
\newcommand{\topic}[1]{\par\addvspace{9pt}\Needspace{10\baselineskip}{\large\bfseries\color{teal}#1}\par\nobreak}
\newcommand{\studentbreak}{\ifsolutions\else\newpage\fi}
\pagestyle{fancy}
\fancyhf{}
\fancyhead[L]{\small\color{muted}Week 2 \textbar\ Problem Set 2}
\fancyhead[R]{\small\color{muted}Patterns, Proofs, and Primes}
\fancyfoot[L]{\footnotesize\color{muted}\edition}
\fancyfoot[R]{\small\color{teal}\thepage}
\renewcommand{\headrulewidth}{0.3pt}
\NewDocumentEnvironment{solution}{+b}{%
  \ifsolutions
  \begin{tcolorbox}[enhanced,breakable,colback=mist!45,colframe=teal!45,
    boxrule=0.5pt,arc=1mm,left=3mm,right=3mm,top=2mm,bottom=2mm,
    before skip=7pt,after skip=7pt,
    title={Solution},title after break={Solution (continued)},
    colbacktitle=mist,coltitle=teal,fonttitle=\small\bfseries]
  \small #1
  \end{tcolorbox}
  \fi
}{}

\begin{document}
{\LARGE\bfseries\color{navy}Problem Set 2}\hfill
{\small\color{teal}\edition}\\[4pt]
{\large Euclid, continued fractions, and prime factorization}\\[4pt]
{\small\color{muted}Week 2 \textbar\ 100 points + 10 optional challenge points}

\ifsolutions
\medskip
\textbf{Teaching note.} Estimated working time is 2--3 hours, excluding the
challenge. The continued fractions
require only finite nesting and exact arithmetic. Subpart scores appear in
parentheses. Award proof credit for the reasoning, including a valid alternative
argument; avoid repeatedly penalizing a carried arithmetic error.
\else
\medskip
\textbf{Name:}\ \rule{0.46\linewidth}{0.4pt}\hfill
\textbf{Due:}\ \rule{0.20\linewidth}{0.4pt}
\fi

\medskip
\textbf{Instructions.} Show your divisions and use exact fractions. Justify
every proof or counterexample. You may use the division algorithm, B\'ezout's
identity, and Euclid's lemma. You may use the fundamental theorem of arithmetic
except when proving it in Problem 4.

\textbf{Notation.} A finite simple continued fraction is
\[
  [a_0;a_1,\ldots,a_n]
  =a_0+\cfrac{1}{a_1+\cfrac{1}{\ddots+\cfrac{1}{a_n}}},
  \qquad [a_0]=a_0.
\]
Here $a_0$ is an integer and $a_i$ is a positive integer for $i\ge1$.
Unless otherwise requested, use $a_n\ge2$ when $n\ge1$.

\topic{I. Euclid and continued fractions}

\problem{One calculation, two readings}{14}
\begin{enumerate}
  \item Use the Euclidean algorithm to compute $\gcd(301,133)$. \partpts{5}
  \item Work backward to find integers $x,y$ such that
  $301x+133y=\gcd(301,133)$. Check your answer by substitution. \partpts{4}
  \item Use the quotients from (a) to express $301/133$ as a simple continued
  fraction. Reduce $301/133$ to lowest terms and explain why this does not
  change its continued-fraction expansion. \partpts{5}
\end{enumerate}
\begin{solution}
\begin{enumerate}
  \item The divisions are
  \[
  \begin{aligned}
  301&=2\cdot133+35,&133&=3\cdot35+28,\\
  35&=1\cdot28+7,&28&=4\cdot7+0.
  \end{aligned}
  \]
  Thus the gcd is $7$.
  \item Reverse the divisions:
  \[
  7=35-28=35-(133-3\cdot35)
    =4(301-2\cdot133)-133=4\cdot301-9\cdot133.
  \]
  Hence $(x,y)=(4,-9)$; indeed $1204-1197=7$.
  \item Dividing each line in (a) by its divisor gives the nested expression
  \[
  \frac{301}{133}=2+\cfrac1{3+\cfrac1{1+\cfrac14}}
    =[2;3,1,4]=\frac{43}{19}.
  \]
  Dividing every entry in the Euclidean divisions by $7$ leaves the
  quotients $2,3,1,4$ unchanged. Equivalently, the continued-fraction
  procedure acts on the value of the rational number, which is unchanged
  by reduction.
\end{enumerate}
\end{solution}

\problem{Unfolding a continued fraction}{12}
\begin{enumerate}
  \item Evaluate $[2;3,4]$ and $[0;2,1,3]$ as fractions in lowest terms. \partpts{4}
  \item For $n\ge1$ and $a_n\ge2$, prove the identity
  \[
  [a_0;a_1,\ldots,a_n]
    =[a_0;a_1,\ldots,a_{n-1},a_n-1,1].
  \]
  Explain the equality by looking at the end of the nested fraction. \partpts{4}
  \item Give a second expansion, ending in $1$, for each number in (a). \partpts{4}
\end{enumerate}
\begin{solution}
\begin{enumerate}
  \item Working from the inside outward,
  \[
  [2;3,4]=2+\frac{1}{3+1/4}
   =2+\frac4{13}=\frac{30}{13},\qquad
  [0;2,1,3]=\frac{1}{2+1/(1+1/3)}
   =\frac{1}{2+3/4}=\frac4{11}.
  \]
  The displayed fractions are reduced: $30=2\cdot13+4$, $13=3\cdot4+1$,
  and $\gcd(4,11)=1$.
  \item The final tail has the same value because
  $[a_n-1;1]=(a_n-1)+1/1=a_n$. Substituting equal tails into the same
  preceding nesting gives equal fractions. The assumption $a_n\ge2$
  ensures that the new entry $a_n-1$ is positive.
  \item The second expansions are
  $30/13=[2;3,3,1]$ and $4/11=[0;2,1,2,1]$.
\end{enumerate}
\end{solution}

\problem{Bounding a square root with fractions}{12}
Consider the four truncations
\[
  c_0=[1],\qquad c_1=[1;2],\qquad c_2=[1;2,2],\qquad c_3=[1;2,2,2].
\]
Such truncations are called \emph{convergents}.
\begin{enumerate}
  \item Compute the four fractions in lowest terms. \partpts{4}
  \item Decide which are below $\sqrt2$ and which are above it by comparing
  their squares with $2$. Do not use decimal approximations. \partpts{4}
  \item Using these values, give the narrowest interval you can with a lower
  endpoint below $\sqrt2$ and an upper endpoint above it. Find its exact
  width. \partpts{4}
\end{enumerate}
\begin{solution}
\begin{enumerate}
  \item Direct nesting gives $c_0=1$, $c_1=3/2$, $c_2=7/5$, and $c_3=17/12$.
  \item For a positive fraction $p/q$, the sign of $(p/q)^2-2$ is the
  sign of $p^2-2q^2$:
  \begin{center}
  \begin{tabular}{@{}cccr@{}}
  \toprule
  Truncation & $p/q$ & $p^2-2q^2$ & Comparison \\
  \midrule
  $c_0$ & $1/1$ & $-1$ & $c_0<\sqrt2$ \\
  $c_1$ & $3/2$ & $1$ & $c_1>\sqrt2$ \\
  $c_2$ & $7/5$ & $-1$ & $c_2<\sqrt2$ \\
  $c_3$ & $17/12$ & $1$ & $c_3>\sqrt2$ \\
  \bottomrule
  \end{tabular}
  \end{center}
  Squaring preserves inequalities between positive numbers, so these
  comparisons prove the asserted bounds.
  \item The larger lower bound is $7/5$, and the smaller upper bound is
  $17/12$. Hence
  \[
  \frac75<\sqrt2<\frac{17}{12},\qquad
  \frac{17}{12}-\frac75=\frac{85-84}{60}=\frac1{60}.
  \]
\end{enumerate}
\end{solution}

\studentbreak
\topic{II. Unique factorization and its consequences}

\problem{Why factorization exists and is unique}{14}
Prove the fundamental theorem of arithmetic in the following steps.
\begin{enumerate}
  \item Use strong induction to prove that every integer $n>1$ is a product
  of primes. Treat the prime and composite cases separately. \partpts{6}
  \item Suppose $p_1\cdots p_r=q_1\cdots q_s$, where all factors are prime.
  Use Euclid's lemma to match and cancel prime factors. Explain why this
  process proves uniqueness up to order, including why one side cannot have
  factors left over. \partpts{6}
  \item Why would allowing $1$ as a prime destroy this uniqueness statement? \partpts{2}
\end{enumerate}
\begin{solution}
\begin{enumerate}
  \item The base case $n=2$ is prime. Assume all integers $m$ with
  $2\le m<n$ are products of primes. If $n$ is prime, it already is such
  a product. If $n$ is composite, then $n=ab$ for integers $1<a,b<n$.
  By the inductive hypothesis, each of $a,b$ is a product of primes;
  multiplying those products gives a prime factorization of $n$.
  \item Since $p_1$ divides $q_1\cdots q_s$, repeated use of Euclid's lemma
  shows that $p_1\mid q_j$ for some $j$. Because $q_j$ is prime and
  $p_1>1$, it follows that $p_1=q_j$. Reorder the $q$'s and cancel this
  factor. Repeat the same argument with the shorter products. The number
  of factors decreases at every step, so the process terminates. If only
  one side had remaining factors, their product would be greater than $1$,
  while the empty product on the other side would equal $1$.
  This is impossible. Thus $r=s$ and the factors agree up to order.
  \item Inserting arbitrarily many $1$'s would give, for example,
  $6=2\cdot3=1\cdot2\cdot3=1\cdot1\cdot2\cdot3$. These lists differ
  in length and cannot be obtained from each other by reordering.
\end{enumerate}
\end{solution}

\problem{Completing a square or a cube}{14}
Let $N=5400$.
\begin{enumerate}
  \item Find the prime factorization of $N$. \partpts{4}
  \item Find the least positive integer $k$ for which $Nk$ is a perfect
  square. Give the square root of $Nk$ and prove that your $k$ is least. \partpts{5}
  \item Find the least positive integer $\ell$ for which $N\ell$ is a
  perfect cube. Give the cube root of $N\ell$ and prove minimality. \partpts{5}
\end{enumerate}
\begin{solution}
\begin{enumerate}
  \item $5400=54\cdot100=2^3\cdot3^3\cdot5^2$.
  \item A square has an even exponent at every prime. The exponents of $2$
  and $3$ in $N$ are odd, so any working multiplier must contain at least
  one factor of $2$ and one of $3$. Thus $6\mid k$ and $k\ge6$.
  Taking $k=6$ works:
  \[
  5400\cdot6=2^4\cdot3^4\cdot5^2
    =(2^2\cdot3^2\cdot5)^2=180^2=32400.
  \]
  \item In a cube, every prime exponent is divisible by $3$.
  The exponents of $2$ and $3$ already satisfy this; the exponent $2$ of
  $5$ needs at least one additional factor of $5$. Therefore every
  working $\ell$ is divisible by $5$ and is at least $5$. This is attained:
  \[
  5400\cdot5=2^3\cdot3^3\cdot5^3=30^3=27000.
  \]
  The exponent tests in both parts follow from uniqueness of prime
  factorization, since taking a square or cube doubles or triples each exponent.
\end{enumerate}
\end{solution}

\problem{True, false, or missing a hypothesis?}{12}
All variables below are positive integers.
\begin{enumerate}
  \item Prove or disprove: if $a^2\mid b^2$, then $a\mid b$. \partpts{3}
  \item Prove or disprove: if $a\mid b^2$, then $a\mid b$. \partpts{3}
  \item Give a counterexample to the claim that $a\mid c$ and $b\mid c$
  imply $ab\mid c$. Then prove that the claim becomes true when
  $\gcd(a,b)=1$. \partpts{6}
\end{enumerate}
\begin{solution}
\begin{enumerate}
  \item True. For any prime $p$, let $\alpha$ be its exponent in $a$
  and $\beta$ its exponent in $b$, allowing zero exponents. Divisibility
  $a^2\mid b^2$ gives $2\alpha\le2\beta$, hence $\alpha\le\beta$.
  This holds for every prime, so $a\mid b$ by unique factorization.
  \item False. Take $a=4$ and $b=2$: $4\mid2^2$, but $4\nmid2$.
  A prime dividing a square must divide its base, but a composite divisor
  can require more copies of a prime than the base contains.
  \item Take $a=4$, $b=6$, and $c=12$. Both $4$ and $6$ divide $12$,
  whereas $24$ does not. Now assume $\gcd(a,b)=1$. No prime occurs in
  both $a$ and $b$. At any prime, its exponent in $ab$ is therefore its
  exponent in one of the factors (or zero). Since each factor divides
  $c$, that exponent is at most the exponent in $c$. Hence $ab\mid c$.
  Alternatively, write $c=au$ and use B\'ezout for $(a,b)$ to deduce
  $b\mid u$, yielding $c=abv$.
\end{enumerate}
\end{solution}

\studentbreak
\topic{III. Finding primes and knowing when to stop}

\problem{The sieve of Eratosthenes}{12}
\begin{enumerate}
  \item Use the grid to sieve the integers from $2$ through $120$. Mark $1$
  as neither prime nor composite. Process primes in increasing order,
  crossing out their multiples beginning at their squares. List all primes
  that remain and give their total number. \partpts{8}
  \item The numbers $91$ and $119$ survive the passes for $2,3,5$.
  Explain why they are nevertheless composite and identify the pass that
  removes them. \partpts{4}
\end{enumerate}

\ifsolutions\else
\begin{center}
\small
\setlength{\tabcolsep}{0pt}
\renewcommand{\arraystretch}{1.32}
\begin{tabular}{|*{10}{>{\centering\arraybackslash}p{0.088\linewidth}|}}
\hline
1&2&3&4&5&6&7&8&9&10\\\hline
11&12&13&14&15&16&17&18&19&20\\\hline
21&22&23&24&25&26&27&28&29&30\\\hline
31&32&33&34&35&36&37&38&39&40\\\hline
41&42&43&44&45&46&47&48&49&50\\\hline
51&52&53&54&55&56&57&58&59&60\\\hline
61&62&63&64&65&66&67&68&69&70\\\hline
71&72&73&74&75&76&77&78&79&80\\\hline
81&82&83&84&85&86&87&88&89&90\\\hline
91&92&93&94&95&96&97&98&99&100\\\hline
101&102&103&104&105&106&107&108&109&110\\\hline
111&112&113&114&115&116&117&118&119&120\\\hline
\end{tabular}
\end{center}
\fi
\begin{solution}
\begin{enumerate}
  \item Since $10^2<120<11^2$, the required sieving primes are $2,3,5,7$.
  Starting at $4,9,25,49$, respectively, and crossing out multiples leaves
  the following $30$ primes:
  \[
  \begin{gathered}
  2,3,5,7,11,13,17,19,23,29,\\
  31,37,41,43,47,53,59,61,67,71,\\
  73,79,83,89,97,101,103,107,109,113.
  \end{gathered}
  \]
  For checking a student's grid, the numbers of composites first removed
  in the four passes are $59,19,7,4$, totaling $89$.
  Thus $1+89+30=120$, including the separate entry for $1$.
  \item $91=7\cdot13$ and $119=7\cdot17$. Both are removed by the pass
  for $7$. Failure to have a factor among $2,3,5$ does not rule out a
  factor of $7$.
\end{enumerate}
\end{solution}

\ifsolutions\newpage\fi
\problem{Why the sieve works}{10}
\begin{enumerate}
  \item Prove that every composite integer $n>1$ has a prime divisor at
  most $\sqrt n$. Deduce that sieving by all primes $p\le\sqrt N$ removes
  every composite integer at most $N$. \partpts{4}
  \item Explain why the pass for a prime $p$ may begin at $p^2$:
  every multiple $mp$ with $2\le m<p$ has already been removed.
  Your argument must also cover composite $m$. \partpts{4}
  \item Would the sieve always work if we used only primes
  $p<\sqrt N$? Give a counterexample showing why the equality matters. \partpts{2}
\end{enumerate}
\begin{solution}
\begin{enumerate}
  \item Write $n=ab$ with $1<a\le b$. If $a>\sqrt n$, then also
  $b>\sqrt n$, so $ab>n$, a contradiction. Therefore $a\le\sqrt n$.
  Existence of prime factorization gives a prime divisor $p$ of $a$;
  then $p\mid n$ and $p\le a\le\sqrt n$. For $n\le N$ this yields
  $p\le\sqrt N$, so a processed prime removes $n$. No prime is removed,
  because every crossed-out number has a proper prime divisor.
  \item Every integer $m\ge2$ has a prime divisor $q\le m$.
  When $m<p$, this gives $q<p$. Thus $q\mid mp$, and the earlier
  pass for $q$ already removed $mp$. Even when the earlier pass began
  at $q^2$, it reaches $mp$ because $m\ge q$ and $p>q$ give
  $mp>q^2$. This reasoning does not assume that $m$ itself is prime.
  \item Take $N=49$. The primes strictly below $\sqrt{49}=7$ are
  $2,3,5$; none divides $49$. The composite $49=7^2$ survives unless
  the pass for $7$ is included. More generally, $N=p^2$ for any prime
  $p$ illustrates the same issue.
\end{enumerate}
\end{solution}

\par\addvspace{12pt}\Needspace{7\baselineskip}
{\bfseries\color{teal}Optional challenge: a square equals a cube}
\pts{10 extra}

Let $a,b$ be positive integers satisfying $a^2=b^3$. Prove that there is a
positive integer $t$ such that $a=t^3$ and $b=t^2$.
\emph{Hint:} compare the exponents of a fixed prime on the two sides.
\begin{solution}
Let the exponents of a prime $p$ in $a$ and $b$ be $\alpha_p$ and
$\beta_p$. Uniqueness of prime factorization gives
$2\alpha_p=3\beta_p$. Since $3\mid2\alpha_p$ and $3\nmid2$,
Euclid's lemma gives $3\mid\alpha_p$. Write $\alpha_p=3\gamma_p$;
the equality then gives $\beta_p=2\gamma_p$.
Set $t=\prod_p p^{\gamma_p}$, a finite product. Its cube is $a$ and its
square is $b$. If $a=b=1$, take $t=1$ (the empty product).
Conversely, any $a=t^3$, $b=t^2$ satisfies $a^2=b^3=t^6$.

\textbf{Suggested credit:} exponent equation (3), divisibility of the
exponents (3), construction and verification of $t$ (4).
\end{solution}

\end{document}
